\chapter{Sonuç ve Öneriler}
\label{ch:sonuc-oneriler}

\section{Sonuç}

Bu tez kapsamında; taksi gibi araç içi ortamlarda sürücü güvenliğini
artırmaya yönelik, uçtan uca çalışan bir gömülü sistem prototipi
tasarlanmış ve gerçekleştirilmiştir. Sistem; bulut tabanlı yüz tanıma
servisi, kullanıcı kaydı / oturum API'si, ESP32-CAM tabanlı kamera +
Wi-Fi + BLE köprüsü, STM32 tabanlı olay yöneticisi ve iPhone üzerinde
SwiftUI ile yazılmış sürücü uygulamasından oluşmaktadır. Tasarım, üç
temel ilkeye sadık kalmıştır:

\begin{enumerate}[label=(\arabic*),itemsep=2pt]
    \item Sürücü etkileşiminin tek bir butona indirgenmesi,
    \item Görsel/işitsel geri bildirimle dikkat dağıtmayan bir kullanıcı
          deneyimi,
    \item Donanım eki maliyetinin 500~TL altında tutulması.
\end{enumerate}

Bulut tarafında açık kaynak InsightFace kütüphanesinin (buffalo\_l, ArcFace
R100) deploy edilmesi, modelin ve eşik değerinin ekip kontrolünde
kalmasını; bu da KVKK uyumlu bir tasarımı mümkün kılmıştır. Çoklu kare
burst yaklaşımı (10 kare $\times$ 5 FPS), hem doğruluk hem pasif canlılık
açısından tek-kare yaklaşıma kıyasla anlamlı kazanç sağlamıştır. STM32
tarafında HAL'dan soyutlanmış bir durum makinesi modülü yazılması, host
PC üzerinde 12 senaryolu birim testin geçirilmesini ve donanım üzerinde
flash--flaş döngülerinin minimuma indirilmesini sağlamıştır.

Sergi öncesi yapılan tezgah testlerinde:

\begin{itemize}[itemsep=1pt]
    \item EC2 sunucusunun health/check, burst smoke testi ve \texttt{/auth/...}
          mutluluk yolu + hata senaryoları geçmiştir;
    \item ESP32-CAM Wi-Fi STA modunda hotspot'a bağlanıp sunucuya 10
          karelik burst gönderebilmektedir;
    \item ESP32-CAM aynı zamanda BLE çevresel olarak \texttt{TaxiGuard}
          adıyla yayınlanmakta, iPhone uygulaması otomatik olarak bağlanıp
          \texttt{FFE1} karakteristiğinin notify aboneliğini açmaktadır;
    \item STM32 CubeMX projesi açılmış, build geçmiş ve kart üzerinde
          HSI$\times$PLL 216~MHz saat üretimi doğrulanmıştır;
    \item iPhone uygulaması Login/Register akışından geçirilmiş, MATCH
          satırı geldiğinde \texttt{tel://}<numara> dialer'ı otomatik
          olarak açılmaktadır.
\end{itemize}

Sergi haftası içinde tamamlanacak adımlar: küçültülmüş bir test seti
üzerinde FAR/FRR ölçümlerinin alınması, sonuçların tezin
\cref{ch:sonuclar}'üne işlenmesi ve demo öncesi \texttt{emergencyNumber}
değerinin 155'e geri çekilmesi.

\section{Akademik ve Pratik Katkılar}

Çalışmanın özgün katkıları şu şekilde özetlenebilir:

\begin{itemize}[itemsep=2pt]
    \item Açık kaynak yüz tanıma motorunun (InsightFace), küçük ölçekli
          gömülü donanım (ESP32-CAM + STM32) ile entegre edilerek bir
          taksi senaryosuna uyarlanması;
    \item Pahalı GSM modülü yerine telefonun BLE ağ geçidi olarak
          kullanılmasıyla, toplam donanım maliyetinin 500~TL'nin altında
          tutulması;
    \item Çoklu kare embedding standart sapmasının, ek model yükü
          olmadan \emph{pasif canlılık} sinyali olarak kullanılması;
    \item HAL'dan soyutlanmış, host-side birim test edilebilen taşınabilir
          bir STM32 durum makinesi modülü;
    \item Otomatik FAR/FRR ölçüm aracı ile, tezin tablo ve grafiklerinin
          tekrarlanabilir biçimde üretilmesi.
\end{itemize}

\section{Geliştirme Önerileri}

\subsection{Faz 2: HTTPS, ESP-CAM Kimlik Doğrulama ve Token Süresi}

Mevcut sürüm (Faz 1) sunucusu HTTP üzerinden çalışmakta; sergi/test
ortamı için yeterlidir ancak üretim için yetersizdir. Faz 2'ye geçiş için:

\begin{itemize}[itemsep=1pt]
    \item EC2 önüne Caddy veya nginx reverse proxy konulması; Let's Encrypt
          ile otomatik sertifika alınması,
    \item ESP32-CAM \texttt{config.h}'da \texttt{USE\_TLS=1} ve ISRG Root~X1
          sertifika gömülmesi,
    \item Sunucu tarafında \texttt{SHARED\_SECRET} aktif edilerek
          ESP32-CAM ve harness'ın \texttt{X-Shared-Secret} başlığı
          taşıması.
    \item Kullanıcı oturum token'larına süre sınırlaması (örnek 30 gün)
          ve yenileme akışı; mevcut sürümde tokenlar yalnızca çıkış
          işleminde geçersizleştirilmektedir.
    \item iPhone uygulamasında \texttt{NSAllowsArbitraryLoads} bayrağının
          kaldırılması, ATS uyumlu HTTPS'in zorunlu kılınması.
\end{itemize}

Bu adımlar tasarım açısından hazırdır; yaklaşık yarım günlük bir geçiş
işidir.

\subsection{Edge Inference (Kısmi Yerel Tanıma)}

Mevcut tasarım, Wi-Fi gerektirir. Wi-Fi yokken sistem güvenli moda geçer
(\cref{sec:safe-mode}) ancak tanıma çalışmaz. Hibrit bir tasarımda:

\begin{itemize}[itemsep=1pt]
    \item ESP32-S3 veya benzeri bir SoC üzerinde quantized ArcFace-MobileNet
          ile kabaca bir ön-eleme yapılabilir;
    \item Yalnızca yüksek-belirsizlik durumlarında bulut sorgulanır.
\end{itemize}

Bu, hem Wi-Fi yokluğunda kısmi işlevsellik hem de bulut maliyetlerinde
düşüş sağlar; ancak ek bir geliştirme yatırımı gerektirir.

\subsection{Filo Entegrasyonu}

Birden fazla aracın merkezi panelde takibi için MQTT veya WebSocket
tabanlı bir olay panosu eklenebilir. Sunucu zaten yapılandırılmış JSON
günlüğü üretmektedir; bu günlüklerin bir MQTT broker'a yayınlanması
büyük bir mimari değişiklik gerektirmemektedir. Mevcut SQLite kullanıcı
tablosu, her sürücüye atanmış plaka bilgisini zaten tuttuğu için
``hangi araçta hangi şoför var'' raporu küçük bir SQL sorgusu kadar
yakındır.

\subsection{Android İstemcisi}
\label{sec:android-gelecek}

Android tarafında aynı sistem, \texttt{Intent.ACTION\_CALL("tel:155")}
ile tek-dokunuşsuz çalışacak şekilde genişletilebilir
\cite{android2024telephony}. Sürücü tarafında yanlış pozitif arama
sorunu, MATCH onayı için kısa bir geri sayım (örneğin 5~s
``İptal etmek için dokunun'') ile azaltılabilir. Bu, iOS sandbox'ı
gibi dialer ekranı açan bir yaklaşıma kıyasla şoförün ekrana
dokunmadan da emniyet hattı ile bağlantı kurabilmesini sağlar.

\subsection{Aktif Canlılık ve Çoklu Modalite}

Mevcut pasif canlılık (cross-frame embedding $\sigma$) etkili olmakla
birlikte, üst düzey saldırılara karşı (3B maske, yüksek çözünürlüklü
ekran) sınırlı kalabilir. Hareket veya kafa-jest tabanlı aktif challenge
(``başınızı sağa çevirin'') eklenmesi, sürücü deneyimini bozmadan
arka koltukta sessizce gerçekleştirilebilir.

\subsection{Yasal ve Etik Kapsam}

Sistemin sahaya çıkarılması için:

\begin{itemize}[itemsep=1pt]
    \item KVKK kapsamında ``açık rıza'' veya ``kanunlarda öngörülmüş
          hâl'' altında çalışacak hukuki çerçevenin netleştirilmesi,
    \item Aranan kişiler veri tabanına girişin yetkili idari/kolluk
          birimlerince yapılacağına dair operasyonel prosedürün
          tanımlanması,
    \item Yolcunun bilgilendirme yükümlülüğünün (örneğin araç içinde
          görünür yazı veya sesli uyarı) ele alınması gerekmektedir.
\end{itemize}

Bu noktalar mühendislik kapsamı dışındadır ancak sistemin gerçek
konuşlandırmaya yol bulması için ön koşuldur.
